\section{Refleksi, Jurnal, dan Portofolio}
\label{sec:ch17-refleksi}

\subsection{Siklus Refleksi}

Refleksi menggunakan siklus \emph{what, so what, now what}:
\begin{enumerate}[label=\arabic*.,leftmargin=*]
  \item \textbf{what}: apa yang terjadi?
  \item \textbf{so what}: mengapa itu penting?
  \item \textbf{now what}: apa yang akan dilakukan berikutnya?
\end{enumerate}

Siklus ini membantu mahasiswa mengubah pengalaman menjadi
pembelajaran yang terarah.

\subsection{Jurnal Mingguan}

Jurnal mingguan memuat:
\begin{itemize}[leftmargin=*]
  \item materi yang dipelajari;
  \item kesulitan yang dihadapi;
  \item strategi yang dicoba;
  \item hasil dan rencana perbaikan.
\end{itemize}

Jurnal tidak perlu panjang. Yang penting adalah konsistensi dan
kejujuran. Jurnal menjadi bukti proses belajar, bukan hanya hasil.

\subsection{Portofolio}

Portofolio memuat bukti kompetensi:
\begin{itemize}[leftmargin=*]
  \item kode proyek dan tes;
  \item laporan praktikum;
  \item diagram dan catatan desain;
  \item refleksi dan evaluasi diri;
  \item presentasi atau demo.
\end{itemize}

Portofolio dinilai berdasarkan kelengkapan, kualitas, dan refleksi.
Bukan jumlah dokumen, tetapi kedalaman bukti.

\begin{table}[H]
\centering
\caption{Rubrik portofolio}
\label{tab:ch17-portofolio}
\begin{tabularx}{\textwidth}{@{}l X r@{}}
\toprule
\textbf{Aspek} & \textbf{Indikator} & \textbf{Bobot} \\
\midrule
Kelengkapan & Semua milestone terwakili & 30\% \\
Kualitas & Kode, tes, dan laporan bermutu & 30\% \\
Refleksi & Analisis mendalam dan jujur & 25\% \\
Presentasi & Portofolio mudah dipahami & 15\% \\
\bottomrule
\end{tabularx}
\end{table}

\subsection{Evaluasi Diri}

Mahasiswa diminta menilai dirinya pada setiap CPMK dengan skala
1--4, disertai bukti. Evaluasi diri dibandingkan dengan penilaian
dosen untuk menemukan kesenjangan persepsi.

\begin{reflectionbox}
Jawab:
\begin{enumerate}[label=\arabic*.,leftmargin=*]
  \item Kompetensi apa yang paling berkembang selama satu semester?
  \item Kompetensi apa yang masih perlu diperkuat?
  \item Bukti apa yang menunjukkan kemajuan tersebut?
  \item Apa rencana konkret untuk semester berikutnya?
\end{enumerate}
\end{reflectionbox}
